% chapters/introduction/introduction.tex
% Texto introductorio del capítulo — aparece antes de la primera sección.
% ============================================================================
%
% GUÍA: 1-2 párrafos que presenten brevemente el contexto general del TFG.
%       No repetir lo que irá en las secciones; sirve como hilo conductor.
%       Mencionar que se trata de un TFG coordinado con Alejandro Soler
%       Sichling (frontend) y que este documento se centra en el backend.
%
% Ejemplo de arranque:
%   La evaluación de exámenes escritos es una tarea central en la actividad
%   docente universitaria. En asignaturas con cientos de alumnos, la
%   corrección manual supone una carga de trabajo considerable...
%   Este Trabajo Fin de Grado forma parte de un proyecto coordinado junto
%   con el TFG de Alejandro Soler Sichling~\cite{soler2026frontend}, que
%   aborda el frontend del sistema...
% ============================================================================

% TODO: Redactar la introducción del capítulo.
